\begin{textsample}{NEG Dim 2 – human – Score -1.00 – sp/sp\_ef\_ciencias\_humanas\_his\_1\_p7.txt}  \label{ex:f2_neg_005}
Na BNCC, um dos principais objetivos do componente curricular é estimular a autonomia de pensamento por intermédio do reconhecimento de diferentes sujeitos, histórias, condutas, modos de ser, agir e pensar sobre o mundo.
Tal percepção estimula o pensamento crítico, pois ajuda a compreender que os indivíduos agem de acordo com a época e o lugar nos quais vivem, o que sintetiza uma ope¬ração fundamental na construção do conhecimento histórico, qual seja, a contextualização.
Rusen ( 2001 ) corro-bora com essa ideia quando afirma que “ a resistência dos homens à per¬da de si e seu esforço de auto-afirmação constituem -se como identidade mediante representações de \textbf{continuidade}, com as quais relacionam as experiências do tempo com as intenções do tempo ” ( p.66 ).

% matched lemmas: continuidade
\end{textsample}
